\section{Related Work}
\label{sec:related}

Our test combines ideas from dataset-inference, model-IP forensics, and
distillation-provenance testing, but changes the null that output-likelihood evidence
is compared against. Dataset inference supplies the statistical template: ownership
or contamination claims should be made through calibrated, many-sample tests rather
than isolated high-likelihood examples
\citep{maini2021datasetinference, maini2024llmdatasetinference,
oren2024testsetcontamination, zhao2025posthocdatasetinference,
dziedzic2022selfsupervised}. TIFS ownership-verification work pushes the same
principle into black-box IP settings through membership fingerprints, backdoor
dataset ownership verification, point-cloud ownership tests, and copyright-evasion
stress cases \citep{liu2022membershipfingerprint, li2023dvbw, wei2025pointncbw,
gao2026ceat2i}. The common lesson is that a forensic null must actually match the
innocent alternative; otherwise distribution shift or task quality becomes a false
ownership signal. We show the analogous failure for distillation provenance: the
base-relative surplus is not a descent statistic until candidate capability is
matched.

Invasive fingerprints and watermarks solve a different problem by pre-injecting a
signal before release. Radioactive data and LLM-text watermarks leave detectable
residuals after training or fine-tuning
\citep{sablayrolles2020radioactivedata, sander2024radioactive}; more recent
anti-distillation and instruction-fingerprinting schemes deliberately amplify or
embed such signals in student behavior
\citep{xu2026antidistillation, yang2026askingback,
xu2024instructionalfingerprinting}. These methods are powerful when the defender
controls generation or training, but they do not apply to an already-released
suspect whose traces were not pre-watermarked. CML is therefore intentionally
post-hoc and non-invasive: it asks what can be supported from suspect outputs plus
authorized candidate likelihood scoring.

The information-forensics literature also sets the review bar for model-IP evidence.
Non-intrusive DNN fingerprints, benign-input fingerprints, model-stealing defenses,
and robust watermarking systems all make the asset, access model, false-accusation
risk, and evasion surface explicit
\citep{zheng2022dnnfingerprint, maho2023fbi, zhang2023cip, zhang2023apmsa,
chen2025queen, gao2024augsteal, li2023linearfunctionality, fei2024wideflat,
zhou2024inversionguided, xu2025ipprotection, yan2025robustnessfingerprint}. Our
asset is narrower: a teacher's reasoning-trace distribution rather than a classifier
checkpoint or private dataset. This makes an honest same-base competitor the critical
false-accusation control, and it is exactly the control that the base-relative
surplus lacks.

Non-invasive model and API fingerprints identify which model or endpoint generated
responses, often with calibrated tests or crafted probes
\citep{pasquini2025llmmap, gao2025modelequalitytesting, gubri2024trap,
jin2024proflingo}. They answer model identity, equality, or family membership, not
whether a suspect was trained from a specified teacher's reasoning traces. Matching
an API to a known family is insufficient for a trace-distillation claim, because an
independently-trained same-base model can resemble the candidate teacher without
descending from it.

The closest prior work studies distillation detection and teacher attribution.
\citet{shi2025distillationdetection} detect distillation with access to the suspect
weights; \citet{wadhwa2025whotaught} perform closed-set teacher identification from
outputs; \citet{liu2025wheredid} use white-box per-action probabilities from the
teacher, original student, and distilled model; and \citet{lee2025quantification}
measure a degree of distillation through behavioral similarity. These works establish
that likelihood agreement can be useful, but they either require white-box evidence,
report closed-set accuracy without an out-of-set false-accusation control, or do not
isolate capability from descent. General black-box provenance tests
\citep{nikolic2025modelprovenance, qiu2026provenanceset} provide the broader
calibrated lineage frame, but their unrelated-model baselines do not remove the
same-base capability/style confound that dominates reasoning-trace distillation.

The contribution here is the capability-matched reference. By comparing the suspect
and a same-base non-distilled reference under the same candidate teacher, CML turns
output likelihood from a resemblance score into a calibrated forensic test. It detects
same-family distillation at low audited false-positive rate where the natural
base-relative surplus fails, while keeping the claim restricted to the declared
candidate set and matched-reference access model.
